\documentclass[14pt,a4paper,oneside]{extbook}
\usepackage[utf8]{vietnam}
\usepackage{amsmath,amssymb,mathrsfs,fancyhdr,enumerate,multirow,makecell}
\usepackage{tkz-euclide,tikz-3dplot,tikz,tkz-tab}
\usetikzlibrary{arrows.meta,calc,intersections,angles,quotes,backgrounds,shapes.geometric,patterns,through,shadows,positioning}
\usepackage{yhmath}
\usepackage[top=0.7cm, bottom=0.7cm, left=1.0cm, right=1.0cm, headheight=16pt, headsep=3mm, footskip=5mm]{geometry}
\usepackage[hidelinks,unicode]{hyperref}
\usepackage{qrcode}
\usepackage[many]{tcolorbox}
\usepackage{fontawesome}

\pagestyle{fancy}
\fancyhf{}
\fancyhead[L]{\footnotesize\bfseries\color{teal!85!black}\faBook~SỔ TAY KIẾN THỨC TOÁN 5}
\fancyhead[R]{\footnotesize\bfseries\color{blue!85!black}Lớp Toán Thầy Lĩnh -- loptoanthaylinh.vn}
\fancyfoot[C]{\footnotesize\bfseries\color{gray!80!black}-- Sổ tay kiến thức: Bài 40. Tìm tỉ số phần trăm của hai số --}
\renewcommand{\headrulewidth}{0.4pt}

\begin{document}

% --- BANNER TIÊU ĐỀ BÀI HỌC ---
\noindent
\begin{tcolorbox}[
  enhanced,
  colback=teal!8!white,
  colframe=teal!80!black,
  arc=3mm,
  boxrule=1.1pt,
  width=\linewidth,
  halign=center,
  top=2.2mm, bottom=2.2mm,
  left=3mm, right=3mm,
  drop fuzzy shadow
]
  {\large\bfseries\color{teal!95!black}\faBookmark\ BÀI 40. TÌM TỈ SỐ PHẦN TRĂM CỦA HAI SỐ}\\[2pt]
  {\scriptsize\color{black!75}\itshape Dạng toán 1 về Tỉ số phần trăm: Quy tắc chia nhẩm, nhân 100 và thêm kí hiệu % trên 1 trang A4}
\end{tcolorbox}

\vspace{3.2mm}
% --- KHỐI 1: QUY TẮC TÍNH CHUẨN ---
\noindent
\begin{tcolorbox}[
  enhanced,
  title={\footnotesize\bfseries\color{white}\faKey\ 1. Quy tắc tìm tỉ số phần trăm của hai số $a$ và $b$},
  colback=blue!2!white,
  colframe=blue!75!black,
  arc=2.5mm, boxrule=0.8pt, top=1.8mm, bottom=1.8mm, left=2.5mm, right=2.5mm
]
  \footnotesize
  Muốn tìm tỉ số phần trăm của hai số $a$ và $b$, ta làm như sau:\\
  $\bullet$ \textbf{Bước 1:} Tìm thương của $a$ và $b$ ($a : b$).\\
  $\bullet$ \textbf{Bước 2:} Nhân thương đó với $100$ và viết thêm kí hiệu \textbf{\%} vào bên phải tích tìm được.\\
  \begin{tcolorbox}[colback=white, colframe=blue!75!black, boxrule=0.6pt, arc=1.5mm, halign=center, top=1.2mm, bottom=1.2mm]
    {\scriptsize\bfseries\color{blue!90!black}$\text{Tỉ số phần trăm của } a \text{ và } b = (a : b \times 100)\%$}
  \end{tcolorbox}
  \textit{VD:} Tìm tỉ số phần trăm của $15$ và $60$:\\
  $15 : 60 = 0,25 \longrightarrow 0,25 \times 100 = 25 \longrightarrow \mathbf{25\%}$.
\end{tcolorbox}

\vspace{3mm}

% --- KHỐI 2: CÁC DẠNG TOÁN THỰC TẾ ---
\noindent
\begin{minipage}[t]{0.52\textwidth}
  \begin{tcolorbox}[
    enhanced,
    title={\footnotesize\bfseries\color{white}\faBriefcase\ 2. Tính tỉ lệ hoàn thành kế hoạch},
    colback=teal!2!white,
    colframe=teal!75!black,
    arc=2.5mm, boxrule=0.8pt, top=1.8mm, bottom=1.8mm, left=2.2mm, right=2.2mm
  ]
    \footnotesize
    Kế hoạch làm $500$ sản phẩm, thực tế làm $550$:\\
    $\bullet$ Đã thực hiện được:\\
    \quad $550 : 500 = 1,1 = 110\%$ kế hoạch.\\
    $\bullet$ Đã vượt mức kế hoạch:\\
    \quad $110\% - 100\% = \mathbf{10\%}$ kế hoạch.
  \end{tcolorbox}
\end{minipage}%
\hfill
\begin{minipage}[t]{0.46\textwidth}
  \begin{tcolorbox}[
    enhanced,
    title={\footnotesize\bfseries\color{white}\faPercent\ 3. Tỉ lệ tăng/giảm phần trăm},
    colback=purple!2!white,
    colframe=purple!70!black,
    arc=2.5mm, boxrule=0.8pt, top=1.8mm, bottom=1.8mm, left=2mm, right=2mm
  ]
    \footnotesize
    Số dân tăng từ $10\,000$ lên $10\,120$ người:\\
    $\triangleright$ Số người tăng thêm: $120$ người.\\
    $\triangleright$ Tỉ lệ tăng dân số:\\
    \quad $(120 : 10\,000) \times 100\% = \mathbf{1,2\%}$.\\
    (Lấy phần tăng thêm chia cho số ban đầu!).
  \end{tcolorbox}
\end{minipage}

\vspace{3mm}

% --- KHỐI 3: SƠ ĐỒ TƯ DUY ---
\noindent
\begin{tcolorbox}[
  enhanced,
  title={\footnotesize\bfseries\color{white}\faSitemap\ Sơ đồ tư duy: Quy trình tìm tỉ số phần trăm},
  colback=yellow!2!white,
  colframe=teal!85!black,
  arc=3mm, boxrule=0.9pt, top=1.8mm, bottom=1.8mm, left=1.5mm, right=1.5mm,
  halign=center
]
  \begin{tikzpicture}[
    scale=0.92,
    every node/.style={transform shape},
    b1/.style={rectangle, rounded corners=4pt, fill=blue!15, draw=blue!70!black, text=blue!90!black, font=\bfseries\scriptsize, inner sep=4pt, align=center},
    b2/.style={rectangle, rounded corners=4pt, fill=teal!15, draw=teal!70!black, text=teal!90!black, font=\bfseries\scriptsize, inner sep=4pt, align=center},
    b3/.style={rectangle, rounded corners=4pt, fill=green!15, draw=green!70!black, text=green!90!black, font=\bfseries\scriptsize, inner sep=4pt, align=center}
  ]
    \node[b1] (n1) at (-5.8, 0) {1. Lấy $a$ chia cho $b$\\Tìm thương dạng STP};
    \node[b2] (n2) at (0, 0) {2. Nhân thương với 100\\(Dịch phẩy phải 2 chữ số)};
    \node[b3] (n3) at (5.8, 0) {3. Viết kí hiệu \%\\vào bên phải kết quả};

    \draw[teal!85!black, thick, -{Stealth[length=2.5mm]}] (n1) -- (n2);
    \draw[teal!85!black, thick, -{Stealth[length=2.5mm]}] (n2) -- (n3);
  \end{tikzpicture}
\end{tcolorbox}

\vspace{3mm}

% --- KHỐI 4: LÀM TRÒN KHI CHIA KHÔNG HẾT & BẪY TRẮC NGHIỆM ---
\noindent
\begin{minipage}[t]{0.54\textwidth}
  \begin{tcolorbox}[
    enhanced,
    title={\footnotesize\bfseries\color{white}\faRepeat\ Khi phép chia thương kéo dài vô tận},
    colback=green!3!white,
    colframe=green!60!black,
    arc=2.5mm, boxrule=0.8pt, top=1.8mm, bottom=1.8mm, left=2.2mm, right=2.2mm
  ]
    \footnotesize
    Nếu thương có nhiều chữ số thập phân:\\
    $\bullet$ Ta chỉ lấy \textbf{4 chữ số ở phần thập phân} của thương rồi làm tròn để khi nhân $100$ được $2$ chữ số thập phân.\\
    \textit{VD:} $19 : 30 = 0,6333\dots \approx 0,6333 = \mathbf{63,33\%}$.
  \end{tcolorbox}
\end{minipage}%
\hfill
\begin{minipage}[t]{0.44\textwidth}
  \begin{tcolorbox}[
    enhanced,
    title={\footnotesize\bfseries\color{white}\faExclamationTriangle\ Mẹo tránh bẫy trắc nghiệm},
    colback=red!2!white,
    colframe=red!70!black,
    arc=2.5mm, boxrule=0.8pt, top=1.8mm, bottom=1.8mm, left=2.2mm, right=2.2mm
  ]
    \footnotesize
    $\bullet$ \textbf{Lỗi trình bày toán học}: Không được viết $15 : 60 = 0,25 \times 100 = 25\%$ (SAI về mặt dấu bằng). Phải viết: $15 : 60 = 0,25 = 25\%$.\\
    $\bullet$ Chú ý số nào chia cho số nào (đọc kĩ đề: tỉ số của nam so với nữ hay so với cả lớp).
  \end{tcolorbox}
\end{minipage}

\end{document}
